\section{Computational Evidence}\label{sec:evidence}
\subsection{Shared restrictions, adaptive certificates, and refresh costs}
The new exact two-period checks reconstruct the optimized release policy, its local box maxima, normalized participation prices, shared-table stationarity, and recursive root plane at nine rational releases, including active-cap transitions. Exhaustive rational KKT-basis enumeration independently recovers the tree optimum. All 81 anchor/query release pairs satisfy the supporting-plane inequality. A separate two-plane shared-table master has value $3/8$. The declared adaptive policy in Proposition~\ref{prop:adaptiveuniform} passes its full-cell feasibility test and has uniform gain certificate $9/1024$. Fourteen deliberately invalid witnesses or certificates are rejected; these finite checks supplement, rather than replace, the proofs.

Table~\ref{tab:r28-scaling} varies decision dimension, service count, parameter dimension, cache size, and active-cap instability. The family is an explicit direct sum of two-period accepted-service blocks with shared child tiers, not a claim that a deep nonrecombining tree has been compressed for free. All generated quantities and audits use rational arithmetic. The same exact full-policy generator and its audit are charged to both pipelines. The fresh comparator is a specialized exact solve and meets a stronger zero-gap criterion within the common target; it is not an artificially expensive generic solver.

The strict tolerance triggers refresh on 83 of 84 queries. All post-governance comparator errors are within the declared target, but the cache pipeline is not faster in these measurements. This is useful adverse evidence: global validity alone does not imply that a cached price is economically tight or computationally beneficial. The full phase accounting and isolated-process memory are reported in EC Section~\ref{ec:r28-evidence}. Historical continuous-state approximation, full-tree comparisons, and unfavorable learning results follow unchanged and are not relabeled as new experiments.
\begin{table}[t]\centering\small
\caption{New exact-arithmetic scaling and cache governance. Each row has 12 queries; the tolerance is $10^{-3}$ per independent block. ``Changes'' is the fraction of block active-cap labels changing between successive queries. Total costs include coefficients, the same full-policy generator and audit, witness repair, cache evaluation, and any refresh. Fresh uses an exact specialized restricted solve. Memory is isolated-process peak resident memory; these are synthetic block-separable models, not dense-QP speed comparisons.}\label{tab:r28-scaling}
\begin{tabular}{rrrrrrrrr}\toprule
$n$ & Services & $d$ & Cache & Changes & Refresh & Cached total & Fresh total & Memory\\
 & & & & (\%) & /12 & (ms) & (ms) & (MiB)\\\midrule
12 & 1 & 1 & 1 & 0.0 & 12 & 1.90 & 1.35 & 16.2 \\
96 & 2 & 4 & 1 & 0.0 & 12 & 15.39 & 11.07 & 16.2 \\
96 & 2 & 4 & 4 & 0.0 & 11 & 20.17 & 11.01 & 16.3 \\
96 & 2 & 4 & 4 & 56.8 & 12 & 20.94 & 11.91 & 16.6 \\
384 & 4 & 8 & 4 & 55.7 & 12 & 81.84 & 45.58 & 16.8 \\
384 & 4 & 8 & 8 & 55.7 & 12 & 105.08 & 44.33 & 17.7 \\
1536 & 4 & 8 & 8 & 55.7 & 12 & 421.12 & 177.34 & 28.3 \\
\bottomrule\end{tabular}
\end{table}


\subsection{Continuous-state guarantees without history enumeration}
Table~\ref{tab:r26-envelopes} reports the inherited continuous-state construction in Section~\ref{sec:continuous}. Bounds cover every inherited tier and feasible root promise, not only evaluation points. In the four-period design, refining from 25 to 289 anchors per date--regime reduces the worst root bound from 0.113253 to 0.019567. The sixteen-period design stores 800 lower anchors rather than constructing a 65,535-node binary history tree; its coarser grid gives a wider 0.193576 bound. These figures report approximation guarantees and storage, not matched-time superiority or a uniform complexity rate. The six designs, complete rational witnesses, and all eleven negative controls are specified in EC Section~\ref{ec:r26}.
\begin{table}[p]\centering
\caption{Continuous promised-state certificates on designed Markov models}\label{tab:r26-envelopes}
\small\begin{tabular}{rrrrrr}\toprule
Dates & Subdivisions & Lower anchors & Upper planes & Uniform gap & Tree nodes\\\midrule
3 & 2 & 54 & 34 & 0.299587 & 7 \\
4 & 4 & 200 & 77 & 0.113253 & 15 \\
4 & 8 & 648 & 153 & 0.051939 & 15 \\
4 & 16 & 2,312 & 322 & 0.019567 & 15 \\
8 & 4 & 400 & 175 & 0.136908 & 255 \\
16 & 4 & 800 & 362 & 0.193576 & 65,535 \\
\bottomrule\end{tabular}
\begin{minipage}{.97\textwidth}\footnotesize Lower anchors and upper planes are totals over all dates and both regimes. The uniform gap is the largest certified root loss over both regimes, all inherited tiers, and all feasible promises, rounded upward to six decimals. Tree nodes give the hypothetical full binary history tree, which is not constructed. These rows are not timing comparisons.\end{minipage}
\end{table}


\subsection{Protocol and common-instance comparator hierarchy}
We use two explicitly different families. The first takes the first 24 nominal contexts of the frozen R24 design: 63 history nodes, two service coordinates, six resource-reward coordinates, all continuation rows, and four shared capacities. A pooled action is an equal \emph{amendment} to the existing heterogeneous tier. The fixed-amendment class is the singleton zero amendment because the exact root payment prevents a nonzero uniform change. Calling this an optimized constant actual tier would be incorrect.

The second family has 63 nodes, one tier, a common baseline of $1/2$, an equiprobable binary public regime, and Markov stage rewards, curvature, and conditional continuation caps. All five classes are feasible and pool actual tiers. The exact root commitment fixes the optimized constant tier at $1/2$ but leaves time-only schedules nontrivial. Every adjacent comparison uses the same primitive instance. We vary reward and switching parameters across 24 seeded contexts; EC Section~\ref{ec:experiments} gives the integer and rational design. Integer probability weights multiply all conditional weights by 32, so its raw value unit differs from the legacy family. We report shares within each family, not a cardinal comparison between the two families.

Table~\ref{tab:hierarchy} reports the optimized increments. Full-history amendments retain substantial value beyond two-period memory on the legacy family. In the public-Markov family most of the gain comes from observing the current regime; the additional memory and full-history gains are smaller and can vanish within the numerical bracket on particular instances. This distinction is informative: the total legacy gain is not presented as a uniquely contractual effect or attributed entirely to learning. Regime-only controls omit inherited tiers and promises even in the Markov family. The exact augmented-state result remains Theorem~\ref{thm:recursion}.

\begin{table}[htbp]
\centering\small
\caption{Value of progressively richer optimized contracts}\label{tab:hierarchy}
\begin{tabular}{lrrrr}
\toprule
& \multicolumn{2}{c}{Legacy amendments} & \multicolumn{2}{c}{Markov service tiers} \\
Expansion & Mean gain & Share (\%) & Mean gain & Share (\%) \\ \midrule
Fixed to time-only & 0.116155 & 7.95 & 0.127557 & 0.49 \\
Time-only to regime & 0.136312 & 9.33 & 26.066328 & 99.34 \\
Regime to two-period memory & 0.224350 & 15.36 & 0.028788 & 0.11 \\
Two-period memory to full history & 0.983916 & 67.36 & 0.016938 & 0.06 \\
\midrule
Total & 1.460733 & 100.00 & 26.239611 & 100.00 \\
\bottomrule
\end{tabular}
\par\smallskip
\begin{minipage}{.98\textwidth}\footnotesize There are 24 common instances per family and five optimized classes per instance. Means use certified lower endpoints; the complete enclosing increment intervals are in the data. Maximum optimization-bracket widths are 8.13e-08 and 2.97e-07, respectively. Shares divide each mean increment by the within-family total. Raw value units differ between families. Current regime omits inherited tier and remaining promise.\end{minipage}
\end{table}


Each optimized class has a rational feasible policy and a separate rational dual upper bound. If these are $[L_j,U_j]$, the certified interval for an increment is
$[\max\{0,L_j-U_{j-1}\},U_j-L_{j-1}]$. We do not interpret a tiny difference between two approximate lower endpoints as a negative true increment. Both the lower-endpoint summaries and these enclosing intervals are retained in the raw record. Exact symbolic tests independently verify the three-node release path, the capacity rent example, and the optimized friction example. Additional rational matrix tests verify the regular-cell release formula. Finite-action promise recursions match exhaustive history enumeration at every attainable root promise; these finite checks supplement the induction proof rather than replace it.

\subsection{Radius refinement and comparator reuse}
Table~\ref{tab:radius} recomputes the true and common-outer time-only comparators on eight fixed legacy contexts and ten radii, including much smaller positive radii than the earlier study. Each hidden coefficient vector is held fixed while its radius varies. The outer-set slack tends toward zero and rises steeply before the earlier smallest positive radius. Thus the coarse positive-radius plateau did not establish a discontinuity. The construction already depended continuously on radius; Proposition~\ref{prop:outer} explains its local behavior. This experiment isolates the outer-set contribution and does not conflate it with objective interpolation over uncertainty vertices.
\begin{table}[htbp]
\centering\small
\caption{The common outer comparator becomes locally tight}\label{tab:radius}
\begin{tabular}{lr}
\toprule
Uncertainty radius & Mean upper bound on outer-set slack \\ \midrule
$0$ & 0.000000003 \\
$1/65536$ & 0.000047365 \\
$1/16384$ & 0.000189415 \\
$1/4096$ & 0.000757072 \\
$1/1024$ & 0.003019002 \\
$1/256$ & 0.011836705 \\
$1/64$ & 0.045064142 \\
$1/16$ & 0.073768497 \\
$1/8$ & 0.074890218 \\
$1/4$ & 0.078860565 \\
\bottomrule
\end{tabular}
\par\smallskip
\begin{minipage}{.98\textwidth}\footnotesize Eight fixed contexts and ten radii, with a directly optimized true and outer comparator at each pair (160 certificates). Entries average the outer certified upper value minus the true certified lower value. At zero the true outer-set slack is exactly zero; the displayed residual is the optimization bracket. Maximum individual bracket width is 3.06e-08. This isolates outer-set loss and is not a calibrated uncertainty experiment.\end{minipage}
\end{table}


For comparator reuse we fix the legacy time-only geometry and switching coefficient, build four certified anchors, and validate 32 nearby reward queries. The rational trace curvature bound is approximately $0.45830$. The largest envelope-minus-certified-lower-value gap is reported in Table~\ref{tab:cache}. The deployed envelope makes no query optimization calls; the 32 validation solves are offline work used to measure tightness. The four initial solves and storage are not free. This is a test of a reusable upper bound and its quality, not a measured end-to-end speedup over the historical robust pipeline.
\begin{table}[htbp]
\centering\small
\caption{Restricted-comparator reuse with fixed geometry}\label{tab:cache}
\begin{tabular}{lr}
\toprule
Quantity & Value \\ \midrule
Certified anchors & 4 \\
Nearby reward queries & 32 \\
Restricted optimizations in deployed query evaluation & 0 \\
Query optimizations used only for offline validation & 32 \\
Curvature trace bound & 0.458298401 \\
Mean envelope slack upper bound & 0.000129719 \\
Maximum envelope slack upper bound & 0.000248488 \\
\bottomrule
\end{tabular}
\par\smallskip
\begin{minipage}{.98\textwidth}\footnotesize A common time-only feasible set and switching coefficient are held fixed. Each query evaluates all four analytical anchor envelopes and takes their minimum. The slack is cached upper value minus independently certified query lower value; it therefore encloses the actual upper-bound loss. Anchor construction, policy acceptance checks, storage, and offline validation are separate costs. No end-to-end timing advantage is inferred.\end{minipage}
\end{table}


\subsection{Correlated comparator certification without query optimization}
Table~\ref{tab:r27-cache} evaluates Theorem~\ref{thm:r27-cache} on the same joint coefficient family as the outer-comparator study, rather than weakening the accepted constraints. Three historical prices per context form a frozen cache. Its first panel contains the 56 non-anchor queries from the finer-radius record. Its second panel contains 24 new parameter vectors outside those inherited one-dimensional rays. A separate true restricted solve, outer solve, and full accepted solve are used only for offline evaluation at each new query. The deployed cache receives none of their optimizers or values.

Across the 24 new queries, the mean certified upper slack decreases from $0.064069$ for the historical outer comparator to $0.016632$ for the three-price cache, with both figures rounded upward. All 24 cache bounds are tighter than their matched outer bounds. The implemented full-class policies have positive cached economic-gain lower bounds in every case, with minimum at least $0.762195$ in this family's objective units. Every optimization bracket is below $6\times10^{-8}$. These are finite, designed comparisons, not a distributional performance or a field-calibration result.

A finite cache is not automatically tight away from its anchors. The one-price column worsens at larger perturbations; adding prices helps without compromising validity. Nor do these calculations establish an end-to-end speed advantage: exact rational arithmetic, all initial solves, the implemented-policy audit, and any requested cache refresh must still be charged. The result is that a new restricted optimization is not logically required to produce the displayed bound. EC Section~\ref{ec:r27-design} specifies all new coefficients, witnesses, source inheritance, and exact checks. Corollary~\ref{cor:r27-cells} separately provides a continuous uncertainty guarantee; the pointwise table is not substituted for that proof.
\begin{table}[htbp]
\centering
\caption{Certified comparator upper slack under correlated parameter changes}
\label{tab:r27-cache}
\begin{tabular}{lrrrr}
\toprule
Radius & Queries & Outer & One price & Three prices\\
\midrule
\multicolumn{5}{l}{\textit{Inherited ray holdouts; no query price reused}}\\
$1/65536$ & 8 & 0.00004737 & 0.00000001 & 0.00000001 \\
$1/16384$ & 8 & 0.00018942 & 0.00000001 & 0.00000001 \\
$1/4096$ & 8 & 0.00075708 & 0.00000011 & 0.00000011 \\
$1/1024$ & 8 & 0.00301901 & 0.00000169 & 0.00000169 \\
$1/256$ & 8 & 0.01183671 & 0.00002694 & 0.00002694 \\
$1/64$ & 8 & 0.04506415 & 0.00043085 & 0.00043085 \\
$1/8$ & 8 & 0.07489022 & 0.02778576 & 0.00711323 \\
\addlinespace
\multicolumn{5}{l}{\textit{New off-ray queries; same frozen cache}}\\
$1/64$ & 8 & 0.04463703 & 0.00073960 & 0.00073960 \\
$1/16$ & 8 & 0.07337423 & 0.01428873 & 0.01428873 \\
$1/8$ & 8 & 0.07419502 & 0.05643124 & 0.03486551 \\
\addlinespace
\bottomrule
\end{tabular}
\par\smallskip
\begin{minipage}{\textwidth}
Entries are means of $U-L_K$, where $L_K$ is an independently checked restricted feasible value. Each decimal is rounded upward. The one-price cache uses the nominal anchor; the three-price cache also uses the inherited radii $1/16$ and $1/4$. No anchor query is included in the first panel. Off-ray queries change coefficient directions, not merely radii. Initial prices and the offline validation optimizations are charged separately; these are certificate-quality comparisons, not matched timings.
\end{minipage}
\end{table}



\subsection{Historical learning and timing evidence}
The R24 matched study charged optimizer labels, fitting, prediction, auditing, polishing, and fallback against lifted continuation. At the strict target both learned routes fell back in every measured case at 63, 127, and 255 nodes. Their mean total times exceeded the lifted comparator in every configuration, with unfavorable paired intervals. The coarse target also produced no observed finite amortization case. Table~\ref{tab:historical} retains the strict-target comparisons; the complete phase-resolved record and all less favorable and more favorable individual cases remain in the code/data archive. These are implementation diagnostics, not a second positive contribution about learned optimization.

Let $C_{\rm off}\geq0$ be a candidate's incremental offline cost, and let $t_C,t_L$ be its and the lifted method's per-query times. The candidate-minus-lifted total cost after $Q_0$ queries is
\begin{equation}\label{eq:amortization}
 C_{\rm off}+Q_0(t_C-t_L)=C_{\rm off}-Q_0\delta,
 \qquad \delta=t_L-t_C.
\end{equation}
For positive offline cost a finite positive crossing exists only when $t_C<t_L$, at $Q_0=C_{\rm off}/(t_L-t_C)$. Strict savings require a query count above this crossing. The analysis code uses this convention, and the displayed equation now matches it. The historical observations have no positive learned crossing under that protocol.
\begin{table}[htbp]
\centering\small
\caption{Preserved strict-target end-to-end timing comparisons}\label{tab:historical}
\begin{tabular}{rrrr}
\toprule
Nodes & Lifted (ms) & Tanh (ms) & RBF (ms) \\ \midrule
63 & 10.34 & 18.84 & 18.91 \\
127 & 96.24 & 112.53 & 113.62 \\
255 & 239.93 & 277.86 & 274.96 \\
\bottomrule
\end{tabular}
\par\smallskip
\begin{minipage}{.98\textwidth}\footnotesize Historical R24 measurements, unchanged. Target $10^{-5}$; 32 held-out contexts and three repetitions per method and size. Both learned routes have 100\% fallback at each size. Their paired saving intervals versus lifted continuation are entirely negative. Offline label/fitting costs are retained separately and cannot turn negative online savings into amortization. These are not newly timed R25 measurements.\end{minipage}
\end{table}


The historical robust study compared different achieved gains. Its roughly 20-millisecond differences between protected learned policies and robust maximin are therefore descriptive cost accounting, not matched-quality dominance. We do not use these three isolated method points as a computational advantage claim. Their complete gain--cost coordinates remain in the archive. The reward-only envelope applies under explicitly fixed geometry, while Theorem~\ref{thm:r27-cache} retains jointly changing coefficients. Neither is retroactively substituted into the old timings; their certificate quality is assessed separately.

